\section{Valuative criterion for universal closedness}
\label{section-valuative-criterion-universal-closedness}

\noindent
In Topology, Section \ref{topology-section-proper} there is a
discussion of proper maps as closed maps of topological spaces
all of whose fibres are quasi-compact, or as maps such that
all base changes are closed maps. Here is the corresponding notion
in algebraic geometry.

\begin{definition}
\label{definition-universally-closed}
A morphism of schemes $f : X \to S$ is said to be
{\it universally closed} if every base change
$f' : X_{S'} \to S'$ is closed.
\end{definition}

\noindent
In fact the adjective ``universally'' is often used in this way.
In other words, given a property $\mathcal{P}$ of morphisms
then we say that ``$X \to S$ is universally $\mathcal{P}$''
if and only if every base change $X_{S'} \to S'$ has $\mathcal{P}$.

\medskip\noindent
Please take a look at Morphisms, Section \ref{morphisms-section-proper} for a
more detailed discussion of the properties of universally closed morphisms.
In this section we restrict the discussion to the relationship between
universal closed morphisms and morphisms satisfying the
existence part of the valuative criterion.

\begin{lemma}
\label{lemma-specializations-lift}
Let $f : X \to S$ be a morphism of schemes.
\begin{enumerate}
\item If $f$ is universally closed then specializations lift
along any base change of $f$, see
Topology, Definition \ref{topology-definition-lift-specializations}.
\item If $f$ is quasi-compact and specializations lift
along any base change of $f$, then $f$ is universally closed.
\end{enumerate}
\end{lemma}

\begin{proof}
Part (1) is a direct consequence of
Topology, Lemma \ref{topology-lemma-closed-open-map-specialization}.
Part (2) follows from
Lemmas \ref{lemma-quasi-compact-closed} and
\ref{lemma-quasi-compact-preserved-base-change}.
\end{proof}

\begin{definition}
\label{definition-valuative-criterion}
Let $f : X \to S$ be a morphism of schemes. We say $f$
{\it satisfies the existence part of the valuative criterion}
if given any commutative solid diagram
$$
\xymatrix{
\Spec(K) \ar[r] \ar[d] & X \ar[d] \\
\Spec(A) \ar[r] \ar@{-->}[ru] & S
}
$$
where $A$ is a valuation ring with field of fractions $K$, the
dotted arrow exists. We say $f$ {\it satisfies the uniqueness
part of the valuative criterion} if there is at most one
dotted arrow given any diagram as above (without requiring
existence of course).
\end{definition}

\noindent
A {\it valuation ring} is a local domain maximal among the relation
of domination in its fraction field, see
Algebra, Definition \ref{algebra-definition-valuation-ring}.
Hence the spectrum of a valuation ring has a unique generic point
$\eta$ and a unique closed point $0$, and of course we have the
specialization $\eta \leadsto 0$.
The significance of valuation rings is that any specialization of
points in any scheme is the image of $\eta \leadsto 0$
under some morphism from the spectrum of some valuation ring.
Here is the precise result.

\begin{lemma}
\label{lemma-points-specialize}
\begin{slogan}
Specializations are witnessed by valuation rings.
\end{slogan}
Let $S$ be a scheme. Let $s' \leadsto s$ be a specialization of points of $S$.
Then
\begin{enumerate}
\item there exists a valuation ring $A$ and a morphism
$f : \Spec(A) \to S$ such that the generic point $\eta$ of
$\Spec(A)$ maps to $s'$ and the special point maps to $s$, and
\item given a field extension $K/\kappa(s')$
we may arrange it so that the extension
$\kappa(\eta)/\kappa(s')$ induced by $f$
is isomorphic to the given extension.
\end{enumerate}
\end{lemma}

\begin{proof}
Let $s' \leadsto s$ be a specialization in $S$, and let
$K/\kappa(s')$ be an extension of fields. By
Lemma \ref{lemma-specialize-points}
and the discussion following
Lemma \ref{lemma-characterize-points}
this leads to ring maps $\mathcal{O}_{S, s} \to \kappa(s') \to K$.
Let $A \subset K$ be any valuation ring whose field of fractions is
$K$ and which dominates the image of $\mathcal{O}_{S, s} \to K$, see
Algebra, Lemma \ref{algebra-lemma-dominate}.
The ring map $\mathcal{O}_{S, s} \to A$ induces the morphism
$f : \Spec(A) \to S$, see
Lemma \ref{lemma-morphism-from-spec-local-ring}.
This morphism has all the desired properties by construction.
\end{proof}

\begin{lemma}
\label{lemma-lift-specializations-valuative}
Let $f : X \to S$ be a morphism of schemes.
The following are equivalent
\begin{enumerate}
\item Specializations lift along any base change of $f$
\item The morphism $f$ satisfies the existence part of the
valuative criterion.
\end{enumerate}
\end{lemma}

\begin{proof}
Assume (1) holds. Let a solid diagram as in
Definition \ref{definition-valuative-criterion} be given.
In order to find the dotted arrow we may replace $X \to S$
by $X_{\Spec(A)} \to \Spec(A)$ since after all
the assumption is stable under base change.
Thus we may assume $S = \Spec(A)$.
Let $x' \in X$ be the image of $\Spec(K) \to X$, so
that we have $\kappa(x') \subset K$, see
Lemma \ref{lemma-characterize-points}.
By assumption there exists a specialization $x' \leadsto x$
in $X$ such that $x$ maps to the closed point of $S = \Spec(A)$.
We get a local ring map $A \to \mathcal{O}_{X, x}$ and a ring
map $\mathcal{O}_{X, x} \to \kappa(x')$, see
Lemma \ref{lemma-specialize-points} and the discussion following
Lemma \ref{lemma-characterize-points}. The composition
$A \to \mathcal{O}_{X, x} \to \kappa(x') \to K$ is the
given injection $A \to K$. Since $A \to \mathcal{O}_{X, x}$
is local, the image of $\mathcal{O}_{X, x} \to K$
dominates $A$ and hence is equal to $A$, by
Algebra, Definition \ref{algebra-definition-valuation-ring}.
Thus we obtain a ring map $\mathcal{O}_{X, x} \to A$ and
hence a morphism $\Spec(A) \to X$
(see Lemma \ref{lemma-morphism-from-spec-local-ring} and
discussion following it). This proves (2).

\medskip\noindent
Conversely, assume (2) holds. It is immediate that
the existence part of the valuative criterion holds for
any base change $X_{S'} \to S'$ of $f$ by considering
the following commutative diagram
$$
\xymatrix{
\Spec(K) \ar[r] \ar[d] & X_{S'} \ar[r] \ar[d] & X \ar[d] \\
\Spec(A) \ar[r] \ar@{-->}[ru] \ar@{-->}[rru] & S' \ar[r] & S
}
$$
Namely, the more horizontal dotted arrow will lead to the
other one by definition of the fibre product. OK, so it clearly
suffices to show that specializations lift along $f$.
Let $s' \leadsto s$ be a specialization in $S$, and let
$x' \in X$ be a point lying over $s'$. Apply
Lemma \ref{lemma-points-specialize}
to $s' \leadsto s$ and the extension of fields
$K = \kappa(x')/\kappa(s')$.
We get a commutative diagram
$$
\xymatrix{
\Spec(K) \ar[rr] \ar[d] & & X \ar[d] \\
\Spec(A) \ar[r] \ar@{-->}[rru] &
\Spec(\mathcal{O}_{S, s}) \ar[r] & S
}
$$
and by condition (2) we get the dotted arrow.
The image $x$ of the closed point of $\Spec(A)$
in $X$ will be a solution to our problem, i.e.,
$x$ is a specialization of $x'$ and maps to $s$.
\end{proof}

\begin{proposition}[Valuative criterion of universal closedness]
\label{proposition-characterize-universally-closed}
Let $f$ be a quasi-compact morphism of schemes.
Then $f$ is universally closed if and only if $f$
satisfies the existence part of the valuative criterion.
\end{proposition}

\begin{proof}
This is a formal consequence of
Lemmas \ref{lemma-specializations-lift} and
\ref{lemma-lift-specializations-valuative} above.
\end{proof}

\begin{example}
\label{example-projective-line-universally-closed}
Let $k$ be a field. Consider the structure morphism
$p : \mathbf{P}^1_k \to \Spec(k)$ of the projective
line over $k$, see Example \ref{example-projective-line}.
Let us use the valuative criterion above to prove that $p$
is universally closed.
By construction $\mathbf{P}^1_k$ is covered by two
affine opens and hence $p$ is quasi-compact.
Let a commutative diagram
$$
\xymatrix{
\Spec(K) \ar[r]_\xi \ar[d] & \mathbf{P}^1_k \ar[d] \\
\Spec(A) \ar[r]^\varphi & \Spec(k)
}
$$
be given, where $A$ is a valuation ring and $K$ is its field
of fractions. Recall that $\mathbf{P}^1_k$ is gotten by glueing
$\Spec(k[x])$ to $\Spec(k[y])$ by glueing
$D(x)$ to $D(y)$ via $x = y^{-1}$ (or more symmetrically $xy = 1$).
To show there is a morphism $\Spec(A) \to \mathbf{P}^1_k$
fitting diagonally into the diagram above we may assume that $\xi$
maps into the open $\Spec(k[x])$ (by symmetry). This gives the following
commutative diagram of rings
$$
\xymatrix{
K & k[x] \ar[l]^{\xi^\sharp} \\
A \ar[u] & k \ar[u] \ar[l]_{\varphi^\sharp}
}
$$
By Algebra, Lemma \ref{algebra-lemma-valuation-ring-x-or-x-inverse}
we see that either
$\xi^\sharp(x) \in A$ or $\xi^\sharp(x)^{-1} \in A$.
In the first case we get a ring map
$$
k[x] \to A,
\ \lambda \mapsto \varphi^\sharp(\lambda),
\  x \mapsto \xi^\sharp(x)
$$
fitting into the diagram of rings above, and we win.
In the second case we see that we get a ring map
$$
k[y] \to A,
\ \lambda \mapsto \varphi^\sharp(\lambda),
\ y \mapsto \xi^\sharp(x)^{-1}.
$$
This gives a morphism
$\Spec(A) \to \Spec(k[y]) \to \mathbf{P}^1_k$
which fits diagonally into the initial commutative diagram of this example
(check omitted).
\end{example}





























